%!TEX root = ../hutbthesis_main.tex
\chapter{主要符号、缩写说明与方法对比表格}

\begin{longtable}[]{@{}
  >{\raggedright\arraybackslash}p{(\columnwidth - 2\tabcolsep) * \real{0.1551}}
  >{\raggedright\arraybackslash}p{(\columnwidth - 2\tabcolsep) * \real{0.8449}}@{}}
\caption{主要符号与缩写说明}\label{tab:appendix_symbols}\\
\toprule\noalign{}
\begin{minipage}[b]{\linewidth}\raggedright
\textbf{缩写/符号}
\end{minipage} & \begin{minipage}[b]{\linewidth}\raggedright
\textbf{含义}
\end{minipage} \\
\midrule\noalign{}
\endfirsthead
\toprule\noalign{}
\begin{minipage}[b]{\linewidth}\raggedright
\textbf{缩写/符号}
\end{minipage} & \begin{minipage}[b]{\linewidth}\raggedright
\textbf{含义}
\end{minipage} \\
\midrule\noalign{}
\endhead
\bottomrule\noalign{}
\endlastfoot
ROS & Robot Operating System，机器人操作系统 \\
Gazebo & 机器人三维物理仿真平台 \\
URDF & Unified Robot Description Format，机器人统一描述格式 \\
A-SMGCS & 高级场面活动引导与控制系统 \\
LiDAR & 激光雷达 \\
\(v\) & 车辆纵向速度 \\
\(\omega\) & 车辆偏航角速度 \\
\(\delta_{f}\) & 前桥转角 \\
\(\delta_{r}\) & 后桥转角 \\
\(e_{x}\) & 纵向误差 \\
\(e_{y}\) & 横向误差 \\
\(e_{\theta}\) & 姿态误差 \\
\end{longtable}

\begin{longtable}[]{@{}
  >{\raggedright\arraybackslash}p{(\columnwidth - 8\tabcolsep) * \real{0.1297}}
  >{\raggedright\arraybackslash}p{(\columnwidth - 8\tabcolsep) * \real{0.2442}}
  >{\raggedright\arraybackslash}p{(\columnwidth - 8\tabcolsep) * \real{0.2060}}
  >{\raggedright\arraybackslash}p{(\columnwidth - 8\tabcolsep) * \real{0.1866}}
  >{\raggedright\arraybackslash}p{(\columnwidth - 8\tabcolsep) * \real{0.2336}}@{}}
\caption{自动牵引对接方法对比}\label{tab:appendix_method_compare}\\
\toprule\noalign{}
\begin{minipage}[b]{\linewidth}\raggedright
\textbf{对比对象}
\end{minipage} & \begin{minipage}[b]{\linewidth}\raggedright
\textbf{主要特点}
\end{minipage} & \begin{minipage}[b]{\linewidth}\raggedright
\textbf{优点}
\end{minipage} & \begin{minipage}[b]{\linewidth}\raggedright
\textbf{局限性}
\end{minipage} & \begin{minipage}[b]{\linewidth}\raggedright
\textbf{与本文方法的关系}
\end{minipage} \\
\midrule\noalign{}
\endfirsthead
\toprule\noalign{}
\begin{minipage}[b]{\linewidth}\raggedright
\textbf{对比对象}
\end{minipage} & \begin{minipage}[b]{\linewidth}\raggedright
\textbf{主要特点}
\end{minipage} & \begin{minipage}[b]{\linewidth}\raggedright
\textbf{优点}
\end{minipage} & \begin{minipage}[b]{\linewidth}\raggedright
\textbf{局限性}
\end{minipage} & \begin{minipage}[b]{\linewidth}\raggedright
\textbf{与本文方法的关系}
\end{minipage} \\
\midrule\noalign{}
\endhead
\bottomrule\noalign{}
\endlastfoot
人工牵引 & 驾驶员和地面人员协同完成牵引作业 & 经验判断灵活，适应真实复杂环境能力强 & 依赖人员经验，标准化和重复精度受人为因素影响 & 本文方法不能直接替代人工，但可为自动化流程提供仿真验证基础 \\
固定路径/开环脚本 & 通过预设速度、转向角和时间完成动作 & 实现简单，适合早期动画和模型验证 & 缺少误差反馈，环境或初始状态变化后适应性差 & 本文在此基础上增加相对位姿误差控制思想 \\
PID控制 & 对速度、角度等单变量进行反馈调节 & 结构简单，工程应用广 & 难以单独处理多变量耦合的精确对接问题 & 可作为底层控制基础，但不足以完整描述对接策略 \\
Pure Pursuit & 基于前视点进行路径跟踪 & 路径跟踪平滑，实现难度适中 & 末端精确姿态对准能力有限 & 适合普通循迹，本文更强调近距离位姿误差修正 \\
Stanley算法 & 基于横向误差和航向误差进行控制 & 适合车辆路径跟踪，横向误差修正能力较好 & 更偏向中高速路径跟踪，低速精细对接针对性不足 & 可作为参考，但本文场景更偏低速对接 \\
MPC控制 & 在约束条件下进行预测优化控制 & 理论效果好，可处理多约束问题 & 建模和调参复杂，实现难度较高 & 后续可作为优化方向，本文阶段暂不采用 \\
本文方法 & ROS-Gazebo仿真、飞机辅助感知、相对位姿误差控制、4WS执行 & 流程完整、结构清晰、实现难度适中，适合毕业设计验证 & 仿真简化较多，真实环境鲁棒性仍需进一步验证 & 作为自动牵引对接系统的前期仿真与算法验证方案 \\
\end{longtable}

\chapter{ROS-Gazebo仿真系统主要命令与自动牵引作业核心代码节选}

\section{主要启动与测试命令}

ROS-Gazebo联合仿真系统的主要启动流程如下所示。

\begin{lstlisting}
cd ~/tug_ws
catkin_make
source devel/setup.bash
roslaunch zongzhuang gazebo.launch
\end{lstlisting}

控制器加载命令如下所示。

\begin{lstlisting}
rosrun controller_manager spawner \
  /zongzhuang/front_left_wheel_velocity_controller \
  /zongzhuang/front_right_wheel_velocity_controller \
  /zongzhuang/back_left_wheel_velocity_controller \
  /zongzhuang/back_right_wheel_velocity_controller
\end{lstlisting}

测试脚本运行命令如下所示。

\begin{lstlisting}
chmod +x test_control.py
rosrun tug_control test_control.py
\end{lstlisting}

\section{目标点生成核心代码}

\begin{lstlisting}[language=Python]
def transform_local_to_world(self, base_pose, local_x, local_y, local_yaw=0.0):
    c = math.cos(base_pose.yaw)
    s = math.sin(base_pose.yaw)
    world_x = base_pose.x + c * local_x - s * local_y
    world_y = base_pose.y + s * local_x + c * local_y
    world_yaw = wrap_to_pi(base_pose.yaw + local_yaw)
    return Pose2D(world_x, world_y, world_yaw)


def generate_task_points(self, aircraft_pose):
    nose_gear_pose = self.transform_local_to_world(
        aircraft_pose,
        self.nose_gear_offset_x,
        self.nose_gear_offset_y,
        0.0
    )

    dock_pose = self.transform_local_to_world(
        nose_gear_pose,
        -self.hitch_offset_x,
        -self.hitch_offset_y,
        0.0
    )

    pre_dock_pose = self.transform_local_to_world(
        dock_pose,
        -self.pre_dock_distance,
        0.0,
        0.0
    )

    pushback_pose = self.transform_local_to_world(
        dock_pose,
        -self.pushback_distance,
        0.0,
        0.0
    )

    retreat_pose = self.transform_local_to_world(
        pushback_pose,
        self.retreat_offset_x,
        self.retreat_offset_y,
        0.0
    )

    return pre_dock_pose, dock_pose, pushback_pose, retreat_pose
\end{lstlisting}

\section{相对位姿误差控制核心代码}

\begin{lstlisting}[language=Python]
def compute_pose_control(self, current_pose, target_pose, max_speed):
    dx = target_pose.x - current_pose.x
    dy = target_pose.y - current_pose.y

    c = math.cos(current_pose.yaw)
    s = math.sin(current_pose.yaw)
    ex = c * dx + s * dy
    ey = -s * dx + c * dy

    distance = math.sqrt(dx * dx + dy * dy)
    target_angle_body = math.atan2(ey, ex)

    direction = 1.0
    alpha = target_angle_body
    if abs(target_angle_body) > math.pi / 2.0:
        direction = -1.0
        alpha = wrap_to_pi(target_angle_body - math.pi)

    yaw_error = wrap_to_pi(target_pose.yaw - current_pose.yaw)
    speed = self.k_rho * distance
    speed = clamp(speed, self.min_motion_speed, max_speed)
    speed *= direction

    if distance > 0.8:
        omega = self.k_alpha * alpha
    else:
        omega = self.k_alpha * alpha + self.k_yaw * yaw_error

    if direction < 0.0:
        omega = -omega

    if abs(speed) < 1e-5:
        curvature = 0.0
    else:
        curvature = omega / speed

    return speed, curvature, distance, abs(yaw_error)
\end{lstlisting}
